% ID: FIS-TER-CAL-001
% Titolo: Confronto tra materiali con diverso calore specifico
% Disciplina: Fisica
% Area: Termologia
% Argomento: Calore specifico
% Sottoargomento: Riscaldamento di materiali diversi
% Classe: Seconda liceo scientifico
% Difficolta: 4
% Tipo: problema_numerico
% Risultato: il rame aumenta di temperatura circa 5,8 volte piu dell'alluminio
% Tempo_stimato: 10 min
% Competenze: confronto quantitativo, proporzionalita inversa, interpretazione fisica
% Prerequisiti: calore specifico, formula Q=mc Delta T
% Tag: termologia, calore_specifico, rame, alluminio, confronto_materiali
% Fonte: esempio_originale
% Autore: Riccardo Carli
% Licenza: CC-BY-SA-4.0

\begin{esercizio}
Due blocchi, uno di rame e uno di alluminio, hanno la stessa massa
$m=0{,}50\,\mathrm{kg}$ e ricevono la stessa quantita di calore
$Q=4500\,\mathrm{J}$. Sapendo che
$c_{\mathrm{rame}}=390\,\mathrm{J/(kg\,K)}$ e
$c_{\mathrm{alluminio}}=900\,\mathrm{J/(kg\,K)}$, determina l'aumento di
temperatura di ciascun blocco e confronta i risultati.
\end{esercizio}

\begin{soluzione}
Usiamo la relazione:
\[
Q=mc\Delta T
\qquad\Rightarrow\qquad
\Delta T=\frac{Q}{mc}.
\]

Per il rame:
\[
\Delta T_{\mathrm{rame}}=
\frac{4500}{0{,}50\cdot 390}
=\frac{4500}{195}
\simeq 23{,}1\,\mathrm{K}.
\]

Per l'alluminio:
\[
\Delta T_{\mathrm{alluminio}}=
\frac{4500}{0{,}50\cdot 900}
=\frac{4500}{450}
=10\,\mathrm{K}.
\]

Il rame aumenta di temperatura piu dell'alluminio perche ha calore specifico
minore: a parita di massa e calore ricevuto, la variazione di temperatura e
inversamente proporzionale al calore specifico.
\end{soluzione}

